% !PRINTED 147
\weberpage{147}
%<<CONT

\weq{25}{\vartheta_{00}\left(0,\frac{\omega}{2}\right)^{2}=\vartheta_{00}^{2}+\vartheta_{10}^{2}=(1+\varkappa)\,\vartheta_{00}^{2}.}

According to \S\,42 we thus obtain a relation between elliptic functions with two different moduli, and if $\lambda$, $L$, $L'$ arise from $\varkappa$, $K$, $K'$ by the interchange $\left(\omega,\dfrac{\omega}{2}\right)$, then by (23), (24), (25)

\weq{26}{\lambda=\frac{2\sqrt{\varkappa}}{1+\varkappa},\quad L=(1+\varkappa)\,K,\quad L'=(1+\varkappa)\,K'}

\weq{27}{\operatorname{sn}\bigl[(1+\varkappa)\,v,\ \lambda\bigr]=\frac{(1+\varkappa)\operatorname{sn}v}{1+\varkappa\operatorname{sn}^{2}v},}

if for the sake of clarity the modulus is also given as second argument under the function sign sn.

For the \gesperrt{Landen} transformation we have
\[
\begin{aligned}
\frac{\vartheta_{11}(2u,\,2\omega)}{\vartheta_{01}(2u,\,2\omega)}&=\frac{\vartheta_{10}(u)\,\vartheta_{11}(u)}{\vartheta_{00}(u)\,\vartheta_{01}(u)},\\[4pt]
\frac{\vartheta_{01}(0,\,2\omega)^{2}}{\vartheta_{00}(0,\,2\omega)^{2}}&=\frac{2\,\vartheta_{00}\vartheta_{01}}{\vartheta_{00}^{2}+\vartheta_{01}^{2}}=\frac{2\sqrt{\varkappa'}}{1+\varkappa'},\\[4pt]
2\,\vartheta_{00}^{2}(0,\,2\omega)&=(1+\varkappa')\,\vartheta_{00}^{2},
\end{aligned}
\]
and from this, if $\lambda$, $\lambda'$, $L$, $L'$ arise by the interchange $(\omega,2\omega)$ from $\varkappa$, $\varkappa'$, $K$, $K'$:

\weq{28}{\begin{aligned}
\lambda'&=\frac{2\sqrt{\varkappa'}}{1+\varkappa'}, \qquad \sqrt{\lambda}=\frac{\varkappa}{1+\varkappa'}, \qquad \lambda=\frac{1-\varkappa'}{1+\varkappa'},\\[4pt]
2L&=(1+\varkappa')\,K, \qquad\qquad\qquad L'=(1+\varkappa')\,K',
\end{aligned}}

\weq{29}{\operatorname{sn}\bigl[(1+\varkappa')\,v,\ \lambda\bigr]=\frac{(1+\varkappa')\operatorname{sn}v\operatorname{cn}v}{\operatorname{dn}v},}

and in (26) to (29) one recognises, by \S\,42, the formulae (2) and (4) of \S\,9.

\begin{center}
\textbf{\S\,45. \gesperrt{The linear transformation of the elliptic functions.}}
\end{center}

The effect of the linear transformation upon the elliptic functions is best surveyed from the representation by the $\sigma$-functions. From (4) or (10), \S\,42, and the formulae (4), (6) of \S\,37 one obtains:
%CONT>>
